% =====================================================================
%  Reproducing the CIC-IDS2017 Thursday-Morning Web-Attack Scenario
%  in an Airbus CyberRange Testbed
%
%  Self-contained Overleaf document. Requires: booktabs, listings,
%  xcolor, hyperref, geometry (all standard on Overleaf).
% =====================================================================
\documentclass[11pt]{article}
\usepackage[margin=1in]{geometry}
\usepackage{booktabs}
\usepackage{listings}
\usepackage{xcolor}
\usepackage{hyperref}
\lstset{basicstyle=\ttfamily\footnotesize,breaklines=true,frame=single,
        backgroundcolor=\color{gray!8},columns=fullflexible,keepspaces=true}

\title{Reproducing the CIC-IDS2017 Thursday-Morning Web-Attack Scenario\\in an Airbus CyberRange Testbed}
\author{}
\date{}

\begin{document}
\maketitle

\begin{abstract}
We reproduce the Thursday-morning ``Web Attacks'' scenario of the widely used
CIC-IDS2017 intrusion-detection dataset~\cite{sharafaldin2018} inside an Airbus
CyberRange environment. Our objective was to follow the original methodology as
closely as the environment allowed---identical victim addressing, the same three
web attacks (Brute Force, XSS, SQL Injection), a benign background profile, and a
mirror-port capture---and to produce a labelled, flow-based dataset in the same
feature format. This document details the complete network topology and its IP
plan, the software installed on each host, the scripts used to generate and
capture traffic, and a thorough comparison of the resulting dataset against the
original CIC file, including a clear discussion of where our dataset matches the
reference and where (and why) it necessarily differs.
\end{abstract}

% =====================================================================
\section{Introduction}
Intrusion-detection research depends on realistic, labelled network datasets.
CIC-IDS2017 is one of the most widely adopted such datasets; its Thursday-morning
segment contains three application-layer web attacks executed against a
vulnerable web server, embedded in benign background traffic. Because the
original data-generation tooling (notably the B-Profile benign agent) is not
publicly available, exact reproduction is not possible; however, the topology,
attack scenarios, capture method and feature format are well documented and can
be faithfully recreated. This work recreates the Thursday-morning scenario in a
controlled CyberRange, documents every implementation choice, and quantifies the
similarity between our dataset and the original.

% =====================================================================
\section{Testbed Topology and IP Plan}
Following the CIC design, the testbed is partitioned into two isolated networks
joined by a firewall: an \emph{Attack-Network} (\texttt{205.174.165.0/24}) hosting
the attacker, and a \emph{Victim-Network} (\texttt{192.168.10.0/24}) hosting the
victim web server, supporting servers and client workstations. A SPAN (mirror)
port on the victim-network switch duplicates all traffic to a dedicated capture
host. Table~\ref{tab:hosts} lists every host and its address.

\begin{table}[h]
\centering
\caption{Complete host inventory and IP addressing used in the experiment.}
\label{tab:hosts}
\begin{tabular}{@{}llll@{}}
\toprule
Network & Host & IP address & Role in this experiment \\
\midrule
Attack  & Kali (\texttt{attacker-73})   & 205.174.165.73 & Attacker (sole attacker used) \\
Attack  & \texttt{attacker-69/70/71}    & 205.174.165.69/70/71 & Present but unused on Thursday \\
\midrule
Victim  & Web server (DVWA)             & 192.168.10.50 & Victim web application \\
        & \quad public address (NAT)    & 205.174.165.68 & Public IP of the victim \\
Victim  & Services hub (\texttt{ubuntu12-server}) & 192.168.10.51 & Benign services (HTTP/FTP/SSH/DNS) \\
Victim  & \texttt{dc-dns}               & 192.168.10.3  & Original DNS host (non-functional) \\
Victim  & \texttt{client-17}            & 192.168.10.17 & Benign client (Linux) \\
Victim  & \texttt{client-19}            & 192.168.10.19 & Benign client (Linux) \\
Victim  & \texttt{client-25}            & 192.168.10.25 & Benign client (Linux) \\
Victim  & \texttt{client-05}            & 192.168.10.5  & Benign client (Windows) \\
Victim  & \texttt{client-08}            & 192.168.10.8  & Benign client (Windows) \\
Victim  & \texttt{client-09}            & 192.168.10.9  & Benign client (Windows) \\
Victim  & \texttt{client-14}            & 192.168.10.14 & Benign client (Windows) \\
Victim  & \texttt{client-15}            & 192.168.10.15 & Benign client (Windows) \\
Victim  & \texttt{client-12}            & 192.168.10.12 & Excluded (in use by a colleague) \\
Victim  & Capture host                  & 192.168.10.200 & SPAN capture (passive) \\
\midrule
Gateway & Firewall (victim side)        & 192.168.10.1   & LAN interface of the pfSense firewall \\
Gateway & Firewall (attack side)        & 205.174.165.80 & WAN interface of the pfSense firewall \\
Gateway & Internet gateway              & 192.168.10.4   & Used only during setup/software install \\
\bottomrule
\end{tabular}
\end{table}

\subsection{Correspondence with the original CIC addressing}
The victim web server is deliberately placed at \texttt{192.168.10.50} and
published through the public address \texttt{205.174.165.68}, exactly matching the
CIC-IDS2017 victim web server. The attacker is a Kali host at
\texttt{205.174.165.73}, again matching CIC. The benign client IPs
(\texttt{.5, .8, .9, .12, .14, .15, .17, .19, .25}) correspond to the CIC client
set. Two hosts in the original testbed could not be used here: the dedicated DNS
server (\texttt{.3}) was non-functional in our environment, and one client
(\texttt{.12}) was reserved by a colleague; these are treated as documented
deviations (Section~\ref{sec:dev}).

% =====================================================================
\section{Why the Victim and the Services Hub Are Separate}
A key design decision is that benign traffic is never directed at the victim.
The victim (\texttt{192.168.10.50}) runs the Damn Vulnerable Web Application
(DVWA) and is targeted \emph{only} by the attacker. All benign user activity is
instead directed at a separate services hub (\texttt{192.168.10.51}) that runs
HTTP, FTP, SSH and DNS. This separation ensures that every flow arriving at the
victim on port~80 is malicious, which yields clean, unambiguous ground truth for
labelling. In our capture, of all benign flows, $22{,}899$ are directed to the
hub, whereas the victim receives only attacker traffic. The original CIC testbed
uses the same principle: benign traffic is served by internal servers while the
attacks are aimed at the web victim, and a switch mirror captures both.

% =====================================================================
\section{Software Deployed on Each Host}

\subsection{Victim web server (\texttt{192.168.10.50})}
The victim runs DVWA on Apache with a MariaDB backend:
\begin{lstlisting}
sudo apt install -y apache2 mariadb-server php php-mysqli php-gd libapache2-mod-php git
cd /var/www/html
sudo git clone https://github.com/digininja/DVWA.git dvwa
sudo cp dvwa/config/config.inc.php.dist dvwa/config/config.inc.php
sudo mysql -e "CREATE DATABASE dvwa; \
  CREATE USER 'dvwa'@'localhost' IDENTIFIED BY 'p@ssw0rd'; \
  GRANT ALL PRIVILEGES ON dvwa.* TO 'dvwa'@'localhost'; FLUSH PRIVILEGES;"
sudo systemctl restart apache2
\end{lstlisting}
DVWA is initialised through \texttt{http://192.168.10.50/dvwa/setup.php} and its
security level is set to \emph{Low}. The address \texttt{192.168.10.50} is added
as an interface alias so the victim matches the CIC IP while preserving the
host's original address.

\subsection{Benign services hub (\texttt{192.168.10.51})}
The hub provides the benign application protocols:
\begin{lstlisting}
sudo apt install -y apache2 vsftpd openssh-server bind9 bind9utils
sudo useradd -m benign
echo 'benign:benign123' | sudo chpasswd
echo "benign test page" | sudo tee /var/www/html/index.html
\end{lstlisting}
After installation the hub listens on HTTP~(80), FTP~(21), SSH~(22) and DNS~(53).

\subsection{Attacker (\texttt{205.174.165.73})}
The Kali host already provides \texttt{patator}, \texttt{sqlmap}, \texttt{curl}
and \texttt{python3}. A small brute-force wordlist is placed at
\texttt{/tmp/pw.txt}.

% =====================================================================
\section{Traffic-Generation and Capture Scripts}
Scripts are transferred to each host with the CyberRange file-upload facility
(destination \texttt{/root} on Linux hosts, \texttt{C:\textbackslash} on Windows)
and executed locally. Table~\ref{tab:scripts} summarises them.

\begin{table}[h]
\centering
\caption{Scripts, their target hosts and their function.}
\label{tab:scripts}
\begin{tabular}{@{}lll@{}}
\toprule
Script & Host(s) & Function \\
\midrule
\texttt{benign\_linux.sh}    & \texttt{.17, .19, .25}          & multi-user HTTP/FTP/SSH/DNS to the hub \\
\texttt{benign\_windows.ps1} & \texttt{.5, .8, .9, .14, .15}   & multi-user HTTP/DNS to the hub \\
\texttt{attack.sh}           & Kali (\texttt{.73})            & timed Brute Force, XSS, SQL Injection \\
\texttt{capture.sh}          & Capture host (\texttt{.200})   & \texttt{tcpdump} on the SPAN interface \\
\texttt{label\_flows.py}     & Analysis workstation           & label flows by attacker IP + time window \\
\bottomrule
\end{tabular}
\end{table}

\subsection{Benign generators: \texttt{benign\_linux.sh} and \texttt{benign\_windows.ps1}}
Each client runs three concurrent simulated users; with eight machines this
yields approximately twenty-four users, comparable to the twenty-five users
modelled by CIC's B-Profile. Each simulated user repeatedly performs a randomly
chosen action---an HTTP request (including GET, POST, form submission and
session cookies), an FTP download, an SSH login, or a DNS query---directed at the
services hub \texttt{192.168.10.51}, pausing a random one-to-eight seconds
between actions to emulate human interaction. The Windows version reproduces the
same behaviour in PowerShell. Because CIC's B-Profile tool is not public, this
generator is a behavioural substitute rather than the original agent.

\subsection{Attack orchestration: \texttt{attack.sh}}
Executed on the Kali host, this script first authenticates to DVWA and sets its
security level to \emph{Low}. It then performs the three web attacks in the CIC
order against the victim's public address \texttt{205.174.165.68} (translated by
the firewall to \texttt{192.168.10.50}):
\begin{enumerate}
  \item \textbf{Web Brute Force} using \texttt{patator} against the login form,
        iterating over the wordlist in \texttt{/tmp/pw.txt};
  \item \textbf{Cross-Site Scripting (XSS)} by submitting reflected and stored
        script payloads to the DVWA XSS endpoints;
  \item \textbf{SQL Injection} using \texttt{sqlmap} against the DVWA SQLi
        endpoint.
\end{enumerate}
The attacks are separated by short idle gaps, and the precise start and stop time
of each phase is written to \texttt{/tmp/attack\_windows.log} for use during
labelling. A representative brute-force invocation is:
\begin{lstlisting}
patator http_fuzz \
  url="http://205.174.165.68/dvwa/vulnerabilities/brute/?username=admin&password=FILE0&Login=Login" \
  method=GET 0=/tmp/pw.txt \
  header="Cookie: PHPSESSID=<session>; security=low" \
  -x ignore:fgrep='incorrect'
\end{lstlisting}

\subsection{Capture: \texttt{capture.sh}}
On the capture host, \texttt{tcpdump} records all mirrored traffic on the SPAN
interface into a single packet capture:
\begin{lstlisting}
sudo tcpdump -i enp9s0 -s 0 -w thursday_morning.pcap
\end{lstlisting}

\subsection{Labelling: \texttt{label\_flows.py}}
Offline, the packet capture is converted to bidirectional flows with a
CICFlowMeter-compatible feature extractor, and each flow is labelled from the
attack-window log: a flow involving the attacker \texttt{205.174.165.73} during a
recorded attack interval receives that attack's label, and all other flows are
labelled \texttt{BENIGN}. Flow direction follows the CIC convention, i.e.\ the
connection initiator (the attacker) is the source and the contacted service (the
victim on port~80) is the destination.

% =====================================================================
\section{How Separate Streams Become One Labelled Dataset}
The attack and benign streams travel to different destinations
(attacker$\rightarrow$\texttt{.50}, benign$\rightarrow$\texttt{.51}) and do not
interact at any host. They are combined only at the switch, whose mirror port
copies every packet---both streams---into a single capture on the capture host.
Mixing therefore occurs at the observation point, not at any end host. After
capture, the single file is separated back into classes during labelling using
the attacker's IP address and the recorded attack time-windows. This mirrors the
original CIC method, in which a switch mirror port ``completely captured all send
and receive traffic to the network.''

\subsection{Collection procedure}
\begin{enumerate}
  \item Synchronise all host clocks (NTP)---required for correct time-window
        labelling.
  \item Start the capture on the capture host.
  \item Start the benign generators on all clients (60-minute duration).
  \item Start the attack orchestration on the Kali host.
  \item After approximately one hour, stop the capture and retain the
        attack-window log.
  \item Extract flow features and apply labels offline.
\end{enumerate}

% =====================================================================
\section{Resulting Dataset}
The capture spans approximately 88 minutes ($165{,}703$ packets) and yields
$28{,}344$ bidirectional flows described by 68 numeric features, using the CIC
column naming. A \emph{bidirectional} flow aggregates both directions of a
conversation---forward (client to server) and backward (server to client)---which
is why the features occur in matched forward/backward pairs. Benign flows
originate predominantly from the Linux clients \texttt{.17}, \texttt{.19} and
\texttt{.25}, with the services hub \texttt{.51} appearing as the server that
replies; attack flows originate from the Kali host \texttt{.73} toward the victim
\texttt{.50}.

% =====================================================================
\section{Comparison with CIC-IDS2017}

\subsection{Overall structure and labels}
\begin{table}[h]
\centering
\caption{Dataset structure and label distribution: original CIC vs.\ ours.}
\label{tab:labels}
\begin{tabular}{@{}lrr@{}}
\toprule
 & CIC & Ours \\
\midrule
Flows (rows)                 & 458{,}968 & 28{,}344 \\
Columns                      & 85 & 70 \\
Shared feature names         & \multicolumn{2}{c}{70} \\
Blank/unlabelled rows        & 288{,}602 & 0 \\
\midrule
BENIGN                       & 168{,}186 & 26{,}322 \\
Web Attack -- Brute Force    & 1{,}507   & 12$^{\dagger}$ \\
Web Attack -- XSS            & 652       & 1{,}290 \\
Web Attack -- SQL Injection  & 21        & 720 \\
\bottomrule
\end{tabular}
\end{table}
\noindent$^{\dagger}$ The low Brute-Force \emph{flow} count is an aggregation
effect: \texttt{patator} uses HTTP keep-alive, so many login attempts share a
small number of TCP connections. Counting distinct attacker connections instead
gives approximately 808 brute-force, 774 XSS and 4--11 SQL-injection connections,
which reproduces CIC's relative ordering (Brute~Force~$\gg$~XSS~$\gg$~SQLi). We
also note that the original CIC file contains $288{,}602$ blank-label rows,
whereas our dataset is fully labelled.

\subsection{Attack endpoints}
\begin{table}[h]
\centering
\caption{Web-attack flow endpoints (all flows use destination port 80/TCP).}
\label{tab:endpoints}
\begin{tabular}{@{}lll@{}}
\toprule
 & CIC & Ours \\
\midrule
Attacker IP (source) & 172.16.0.1 & 205.174.165.73 \\
Victim IP (destination) & 192.168.10.50 & 192.168.10.50 \\
Destination port / protocol & 80 / TCP & 80 / TCP \\
\bottomrule
\end{tabular}
\end{table}
In both datasets the attacker is the flow source and the victim web server is the
destination on port~80. The single systematic difference is the attacker's
recorded address, discussed next.

\subsection{Why some addresses differ}
\label{sec:addrdiff}
The victim web server matches CIC exactly (private \texttt{192.168.10.50}, public
\texttt{205.174.165.68}), because we deliberately assigned these addresses. The
attacker's \emph{recorded} address differs: CIC's Fortinet firewall applied
source network address translation, so the attacker appears in their capture as
the firewall's internal interface \texttt{172.16.0.1}; our pfSense firewall does
not source-translate the attacker, so it appears as its true Kali address
\texttt{205.174.165.73}. This is a property of the firewall configuration, not of
the attack itself, and does not affect the attack's flow characteristics. Two
supporting hosts also differ: the CIC DNS server (\texttt{.3}) was non-functional
in our environment, so DNS was provided by the services hub \texttt{.51}; and one
client (\texttt{.12}) was reserved by a colleague and therefore excluded. These
substitutions change which hosts source the benign traffic but preserve the
overall structure.

\subsection{Feature set and constant features}
Both datasets use CICFlowMeter-style bidirectional flow features; all 70 of our
column names appear in the original CIC file. Our extractor emits 68 numeric
features, a subset of CIC's full set (CIC additionally includes a small number of
rarely used features). Certain features are constant-zero in both datasets
because the underlying TCP events do not occur in web traffic: our dataset has
one such column (\texttt{URG Flag Count}), while the original CIC file has ten
(including \texttt{Bwd PSH Flags}, the forward/backward URG-flag counts, the
CWE-flag count and the six bulk-rate features). Constant features carry no
information and are routinely removed during machine-learning preprocessing.

\subsection{Statistical comparison of benign flows}
\begin{table}[h]
\centering
\caption{Mean of selected features over BENIGN flows.}
\label{tab:benignstats}
\begin{tabular}{@{}lrr@{}}
\toprule
Feature & CIC (mean) & Ours (mean) \\
\midrule
Flow Duration ($\mu$s)   & 12{,}540{,}399 & 88{,}543 \\
Total Fwd Packets        & 15.18 & 2.90 \\
Total Backward Packets   & 18.19 & 1.83 \\
Flow Bytes/s             & 1{,}810{,}402 & 349{,}444 \\
Fwd Packet Length Mean   & 48.67 & 14.15 \\
Bwd Packet Length Mean   & 161.58 & 216.67 \\
ACK Flag Count           & 0.28 & 4.10 \\
\bottomrule
\end{tabular}
\end{table}
The benign distributions differ in shape. CIC's benign profile comprises
twenty-five users over five protocols, including long-lived HTTPS and e-mail
sessions, producing fewer but longer and heavier flows. Our generator issues
many short HTTP, FTP, SSH and DNS requests to an internal hub, producing shorter
flows with fewer packets and lower per-flow duration. Both are legitimate benign
traffic; the difference reflects the different user-behaviour models and protocol
mixes rather than any error.

\subsection{Statistical comparison of attack flows}
\begin{table}[h]
\centering
\caption{Mean of selected features over attack flows.}
\label{tab:attackstats}
\begin{tabular}{@{}lrrr@{}}
\toprule
Feature & Class & CIC (mean) & Ours (mean) \\
\midrule
Total Fwd Packets   & Brute Force & 12.37 & 2.33 \\
Bwd Packet Len Mean & Brute Force & 41.62 & 231.83 \\
Total Fwd Packets   & XSS         & 7.90 & 11.90 \\
PSH Flag Count      & XSS         & 0.96 & 10.25 \\
Total Fwd Packets   & SQL Inj.    & 3.05 & 11.91 \\
Bwd Packet Len Mean & SQL Inj.    & 327.18 & 1{,}349.83 \\
\bottomrule
\end{tabular}
\end{table}
The three attack types are present and correctly labelled in both datasets, and
their relative ordering matches. Absolute feature values differ for three
reasons: the attack tooling is not identical (we use \texttt{patator} and
\texttt{sqlmap}, whereas CIC automated the browser with Selenium); flow
aggregation differs (keep-alive connections group attempts differently); and our
DVWA instance returns different-sized responses (for example, \texttt{sqlmap}
elicits larger database responses, raising the backward packet length for SQL
injection). These differences affect magnitudes, not the identity or ordering of
the attacks.

\subsection{Summary of closeness and differences}
Our dataset is close to the reference in the aspects that define the scenario:
identical victim addressing and port, the same three web attacks executed in the
same order with the same relative volumes, the same bidirectional flow-feature
format with matching column names, and the same mirror-based capture method. It
differs from the reference in aspects that could not be replicated exactly: the
benign generator is a behavioural substitute for the private B-Profile agent
(covering four of the five original protocols); the attacker's recorded IP
reflects a different firewall NAT policy; two supporting hosts were substituted
or excluded; the run is one hour rather than a full morning; and absolute feature
magnitudes differ because of tooling and content differences. None of these
affects the validity of the dataset for intrusion-detection experimentation, and
all are documented below.

% =====================================================================
\section{Documented Deviations from the Original}
\label{sec:dev}
\begin{enumerate}
  \item \textbf{Benign generator.} CIC's B-Profile agent is not publicly
        available; we implemented an equivalent behavioural generator. E-mail
        (SMTP/IMAP) and HTTPS were not provided, so four of the five original
        protocols (HTTP, FTP, SSH, DNS) are covered.
  \item \textbf{Attacker address.} Recorded as \texttt{205.174.165.73} because
        our firewall does not source-translate the attacker, whereas CIC's
        Fortinet firewall produced \texttt{172.16.0.1}. The victim address
        matches CIC.
  \item \textbf{Substituted/excluded hosts.} The original DNS host (\texttt{.3})
        was non-functional, so DNS was consolidated on the services hub
        (\texttt{.51}); one client (\texttt{.12}) was reserved by a colleague and
        excluded.
  \item \textbf{Scale and duration.} Approximately twenty-four users on eight
        machines and a one-hour run, versus twenty-five users and an
        approximately four-hour morning in the original; the attack order, gaps
        and relative volumes were preserved.
  \item \textbf{Client operating systems.} Host \texttt{.25} runs Linux (the
        original used macOS); as flows are keyed by IP, this does not affect the
        flow features.
  \item \textbf{Feature extractor.} A CICFlowMeter-compatible extractor was used;
        the brute-force flow count differs from the original owing to keep-alive
        aggregation, as discussed in Section~\ref{tab:labels}.
\end{enumerate}

\begin{thebibliography}{9}
\bibitem{sharafaldin2018} I.~Sharafaldin, A.~H.~Lashkari, and A.~A.~Ghorbani,
``Toward Generating a New Intrusion Detection Dataset and Intrusion Traffic
Characterization,'' in \emph{Proc. 4th International Conference on Information
Systems Security and Privacy (ICISSP)}, 2018, pp.~108--116.
\end{thebibliography}

\end{document}
